\section{Primary-source ledger and the number-field bridge}
\label{app:sources}

The bibliography is deliberately restricted to the six entries audited
before drafting.  \Cref{tab:source-ledger} records the precise result imported
from each source and, equally important, the instance facts that the source
does not supply.

\begin{longtable}{@{}>{\raggedright\arraybackslash}p{0.13\textwidth}
 >{\raggedright\arraybackslash}p{0.25\textwidth}
 >{\raggedright\arraybackslash}p{0.29\textwidth}
 >{\raggedright\arraybackslash}p{0.25\textwidth}@{}}
\caption{Primary-source locators and claim boundaries.}
\label{tab:source-ledger}\\
\toprule
source & exact locator & use here & not supplied\\
\midrule
\endfirsthead
\toprule
source & exact locator & use here & not supplied\\
\midrule
\endhead
\cite{SwinnertonDyer1993}
& p.~449; \S2; Lemma~1, pp.~451--452; Lemma~2, p.~453; concluding
classification, p.~458
& number-field Picard/Brauer description; invariant-double-six criterion;
complete nonzero 2-primary possibilities \(\Z/2\) and \((\Z/2)^2\)
& the frozen resolvers, fixed fields, orientation square, quartic, or
quaternion representative\\
\addlinespace
\cite{ElsenhansJahnel2010Brauer}
& introduction \S1.3; Lemma~2.5; Remark~3.3 and Notation~3.4;
Theorems~3.5 and 4.4; \S4.3; Definition~4.1; Proposition~5.8;
Theorem~5.11; Fact~6.2, Remark~6.3, Corollary~6.4
& \(U_1/U_3\) containment and restriction; orientation; invariant
double-six class map; cyclic/quaternion precedent
& the degree-\(36\) field or determinant quartic for \cref{eq:cubic};
Corollary~6.4 is not used for a local evaluation\\
\addlinespace
\cite{ElsenhansJahnel2010DoubleSix}
& Theorem~6.6; Algorithm~6.7; Propositions~7.1 and 7.4; Example~9.2
& explicit invariant-double-six descent, tritangent-plane context,
sixer-splitting radicand, and a full-\(U_1\) \([12,15]\) example
& inversion of the present frozen full-\(\WE\) line field\\
\addlinespace
\cite{ElsenhansJahnel2012OrderThree}
& Remark~4.34(iii), printed p.~22
& prior-art boundary: a degree-\(36\) resolver with roots corresponding to
double-sixes
& either resolver polynomial or the Brauer class in this paper\\
\addlinespace
\cite{FarbWolfson2019}
& Theorem~5.6 and Corollary~5.9
& moduli-cover and resolvent-degree context for double-sixes and ordered
sixers
& an arithmetic fixed field for the present surface\\
\addlinespace
\cite{Viray2023}
& \S2.1, p.~10, equation~(24) and the following low-degree exact sequence
& supporting Hochschild--Serre exposition
& any instance computation or automatic Brauer--Manin conclusion\\
\bottomrule
\end{longtable}

\subsection{Why the subgroup theorem applies over
\texorpdfstring{\(\FD\)}{F-D}}

The attaining base is the number field \(\FD\), not \(\Q\).
Swinnerton-Dyer's Picard/Brauer theorem already works over algebraic number
fields \cite{SwinnertonDyer1993}; the Elsenhans--Jahnel containment and
restriction statements concern finite subgroups of \(\WE\), the Picard
lattice, and line configurations \cite{ElsenhansJahnel2010Brauer}.  Thus for
any number field \(k\) they apply to
\[
 H=\im(G_k\to\WE)
\]
without a \(\Q\)-specific step.  For the representative, we construct
\(\FDp/\FD,\cD,f_D\) in the present tower, prove
\(\divv(f_D)=\Norm(\cD)\), calculate its Picard cocycle, and only then use
\cref{eq:brauer-picard-bridge}.

\subsection{Bounded novelty screen}

The 2026-08-15 audit screened exact coefficients and title phrases,
configuration/Brauer queries, and the sources' identifiers.  It supports only:
\begin{quote}
The bounded 2026-08-15 screen did not locate a prior exact computation of the
degree-36 double-six field, its orientation square, and the
determinant-defined quaternion generator for this frozen Yukawa surface.
\end{quote}
This is not an absolute priority claim: the generic constructions remain
prior art.
